\section{Introduction}
\label{sec:intro}

Natural convection between horizontal cylindrical annuli is a canonical problem in thermal engineering, underlying the design of solar collector receivers, thermal energy storage units, electrical machinery with rotating shafts, and gas-insulated transmission cables. Kuehn and Goldstein~\cite{KuehnGoldstein1978} provided the foundational experimental dataset for concentric and eccentric configurations, measuring heat transfer across a wide Rayleigh number range spanning conduction-dominated, laminar convective, and partially turbulent regimes; their ratio of actual to conduction heat transfer remains a standard way to report results in this problem class and is used directly for validation in the present work. Prusa and Yao~\cite{PrusaYao1983} derived a closed-form conformal-mapping solution for pure-conduction heat transfer between eccentric cylinders, and Projahn et al.~\cite{Projahn1981} and Fusegi et al.~\cite{Fusegi1986} independently corroborated the Kuehn and Goldstein benchmark numerically. More recently, Vjatkin et al.~\cite{Vjatkin2024} reported a new experimental study of thermal convection in a rotating horizontal annulus, confirming that this configuration remains an active area of experimental investigation.

Rotation of the inner cylinder introduces a second physical mechanism beyond buoyancy alone: a rotation-induced shear flow at the inner wall that interacts with the buoyancy-driven convective cells. Lee~\cite{Lee1992a,Lee1992b} studied laminar convection between concentric and eccentric heated horizontal cylinders with a rotating inner cylinder, reporting that rotation monotonically suppresses the mean Nusselt number at fixed Rayleigh number, at a rate that depends on eccentricity, a finding the present study independently confirms. Yoo~\cite{Yoo1998} studied the reversed configuration, a rotating outer cylinder with a stationary inner cylinder, and used an exact circular-Couette analytical solution as an isothermal-limit check on their numerical method, a validation strategy adapted here to the present (rotating-inner) configuration. Kaewbumrung and Charoenloedmongkhon~\cite{Kaewbumrung2022} studied three-dimensional turbulent flow in a co-rotating eccentric configuration, providing further evidence that rotation-driven flow-structure changes in this problem class are physically real and numerically resolvable. El-Maghlany and Abo Elazm~\cite{ElMaghlany2016} studied nearly the same base configuration as the present work, an eccentric horizontal annulus with a rotating, heated inner cylinder, adding a nanofluid dimension not used here; this is the closest identified prior computational study of the base physical geometry, and its existence indicates that the geometry itself is well characterized in the literature rather than novel, a distinction addressed explicitly below. Singh et al.~\cite{Singh2024} studied natural convection in an eccentric annulus without rotation across a range of eccentricities and angular positions, and Pranowo et al.~\cite{Pranowo2024} reported, in a related but distinct trapezoidal-cavity geometry, a failure to reach a steady state at $Ra=10^6$, independent evidence that the highest Rayleigh number considered in the present study sits near the edge of the steady laminar regime for this class of problem, a physical limitation of the present study as well.

Computing this configuration's response across a wide parameter space by direct simulation alone is computationally expensive, which motivates data-driven surrogate models: machine learning approximations trained on a finite ensemble of high-fidelity computational fluid dynamics (CFD) solutions that predict thermal performance at new parameter combinations orders of magnitude faster than a direct simulation. Surrogate modeling has seen growing adoption in thermal cavity and annulus flows specifically. Feng et al.~\cite{Feng2024} developed a graph convolution network surrogate for natural convection in an annulus with a variable configuration of internal heat sources; their study varies heat-source number, size, and position rather than the base annulus geometry, includes no rotation or eccentricity, does not use Rayleigh number as a swept model input, and targets the temperature field directly rather than the Nusselt number, making it a materially different problem despite superficial similarity. Zhu et al.~\cite{Zhu2026} recently proposed a physics-informed neural network with geometry-conditioned adapters for transfer learning across eccentricities in a buoyancy-driven enclosure with an inner heated cylinder; based on the published description, this work does not include rotation and is framed around parameter-efficient transfer across geometries rather than a joint sweep over rotation, eccentricity, and Rayleigh number as independent model inputs. Agarwal et al.~\cite{Agarwal2021} demonstrated deep-learning surrogates (feedforward, recurrent, and convolutional-autoencoder architectures) for a two-dimensional mantle convection problem, a different physical domain but a useful methodological precedent for the parameterized spatial-field surrogate developed here. Travnikov et al.~\cite{Travnikov2024} advanced graph neural network architectures specifically for fluid flow and heat transfer surrogate modeling under variable boundary conditions and geometry, and explicitly documented an accuracy degradation for larger computational domains caused by the limited receptive field of standard message-passing graph neural networks, a limitation independently observed for the graph neural network candidate evaluated in the present work. Kokash et al.~\cite{Kokash2024Energies} developed random forest and physics-informed neural network surrogates for a single rotating cylinder in a lid-driven cavity, and the same author group subsequently extended this to two counter-rotating cylinders with a Gaussian process regression surrogate achieving $R^2=0.9999$; both studies establish the methodological lineage that the present work extends to a rotating, eccentric annular geometry.

Despite this growing body of work, no prior study, to the authors' knowledge, has constructed a machine learning surrogate over a joint parameter space of radius ratio, eccentricity, Reynolds number, and Richardson number for a rotating, eccentric, cylinder-in-cylinder annulus. The novelty of the present study therefore rests specifically on this combination, an architecture comparison across a modern gradient-boosted tree ensemble, a Kolmogorov-Arnold Network, and a physics-informed neural network for the scalar Nusselt number target, and a comparison of a per-node neural network, a DeepONet, and a graph neural network for the spatial temperature and velocity field target, rather than on the underlying physical configuration being previously unstudied: as noted above, El-Maghlany and Abo Elazm~\cite{ElMaghlany2016} already established the base geometry as a well-characterized physical system. A further, architecture-independent contribution is methodological: the winning scalar surrogate is validated not only on a held-out split of the training grid but against new simulations mathematically guaranteed to be absent from it, a distinction this study shows is far from academic, since the same-grid figure understates the true error by more than an order of magnitude here, and argues should be adopted more routinely in CFD surrogate-modeling practice.

A methodological choice made early in this study's development deserves explicit mention, since it directly shapes the dataset design below. An initial version of this work parameterized the CFD sweep directly in terms of Rayleigh number and rotation speed, following the convention of much of the classical literature reviewed above. That approach, while internally consistent, conflates two physically distinct effects, buoyancy strength and rotational forcing, in a way that makes it difficult to independently place a given simulation on the natural-to-mixed-to-forced convection spectrum without first computing a derived ratio. The present study instead sweeps viscosity and rotation speed such that the achieved Reynolds number, $Re$, and Richardson number, $Ri=Gr/Re^2$, are the independent, directly-interpretable swept quantities, with the rotation rate itself held fixed at a single reference value rather than swept as a nominally-independent fifth parameter, a choice justified below on physical grounds. This reparameterization is what allows the dataset to span, by construction rather than by post-hoc classification, the full range from buoyancy-dominated ($Ri \gg 1$) through the mixed-convection transition ($Ri \sim 1$) to rotation-dominated ($Ri \ll 1$) flow, with deliberately increased sampling density around $Ri\sim1$, where both mechanisms are comparably important and where the physics is least well captured by either pure limit.

The present study addresses this gap through a systematic surrogate study of a rotating eccentric horizontal annulus. A dataset spanning 35 geometries, 12 Reynolds numbers ($Re=10$ to $1000$), and 13 Richardson numbers ($Ri=0.01$ to $10$, densified near $Ri=1$) is generated and validated against the experimental data of Kuehn and Goldstein~\cite{KuehnGoldstein1978}, the qualitative rotation-effect claims of Lee~\cite{Lee1992a,Lee1992b}, and the independently published closed-form conduction result of Prusa and Yao~\cite{PrusaYao1983}. Three candidate architectures are trained and compared for the scalar Nusselt number surrogate, and three further candidates for the spatial field surrogate, with all candidates reported rather than only the eventual winner. The accuracy claims made for the winning scalar model are validated two ways: against a held-out split of the original simulation grid, the conventional approach, and, separately, against entirely new numerical simulations run at operating conditions mathematically guaranteed to be absent from the training data. The second check reveals a materially higher error than the first, a distinction this study argues should be reported more routinely in surrogate-modeling work of this kind, since the held-out-split accuracy figure alone can substantially overstate real-world predictive reliability.
